\ifdefined\gkver\else\def\gkver{student}\fi
\documentclass[\gkver]{gaokaozhenti}
\begin{document}

\chapter{2005年全国理综Ⅲ}
%% paperType: 理综物理部分

\begin{enumerate}
\item
%% number: 14
%% paperName: 2005年全国理综Ⅲ第 14 题 6 分
%% typeId: 1
%% type: 单选题
%% chapter: 运动和力
%% point: 牛顿第二定律
%% method: 无
%% score: 6
%% degree: 650
%% duplicateId: 0
%% body:
如图所示，一物块位于光滑水平桌面上，用一大小为 $F$、方向如图所示的力去推它，使它以加速度 $a$ 右运动。若保持力的方向不变而增大力的大小，则\xzanswer{A}

\onepicture[w=4.0cm]{figs/14.png}

\fourchoices[answer=A]
{$a$ 变大}
{不变}
{$a$ 变小}
{因为物块的质量未知，故不能确定 $a$ 变化的趋势}

%% answer: A
%% memo:
\memoanswer{
本题原卷及材料中未提供详解，待补充。
}

\item
%% number: 15
%% paperName: 2005年全国理综Ⅲ第 15 题 6 分
%% typeId: 1
%% type: 单选题
%% chapter: 原子和原子核
%% point: 玻尔模型
%% method: 无
%% score: 6
%% degree: 650
%% duplicateId: 0
%% body:
氢原子的能级图如图所示。欲使一处于基态的氢原子释放出一个电子而变成氢离子，该氢原子需要吸收的能量至少是\xzanswer{A}

\onepicture[h=4.5cm]{figs/15.png}

\fourchoices[answer=A]
{$13.60\UeV$}
{$10.20\UeV$}
{$0.54\UeV$}
{$27.20\UeV$}

%% answer: A
%% memo:
\memoanswer{
本题原卷及材料中未提供详解，待补充。
}

\item
%% number: 16
%% paperName: 2005年全国理综Ⅲ第 16 题 6 分
%% typeId: 1
%% type: 单选题
%% chapter: 电磁感应
%% point: 楞次定律
%% method: 无
%% score: 6
%% degree: 650
%% duplicateId: 0
%% body:
如图，闭合线圈上方有一竖直放置的条形磁铁，磁铁的 N 极朝下。当磁铁向下运动时（但未插入线圈内部），\xzanswer{B}

\onepicture[h=4.0cm]{figs/16.png}

\fourchoices[answer=B]
{线圈中感应电流的方向与图中箭头方向相同，磁铁与线圈相互吸引}
{线圈中感应电流的方向与图中箭头方向相同，磁铁与线圈相互排斥}
{线圈中感应电流的方向与图中箭头方向相反，磁铁与线圈相互吸引}
{线圈中感应电流的方向与图中箭头方向相反，磁铁与线圈相互排斥}

%% answer: B
%% memo:
\memoanswer{
本题原卷及材料中未提供详解，待补充。
}

\item
%% number: 17
%% paperName: 2005年全国理综Ⅲ第 17 题 4 分
%% typeId: 1
%% type: 单选题
%% chapter: 电路
%% point: 电容
%% method: 无
%% score: 4
%% degree: 650
%% duplicateId: 0
%% body:
水平放置的平行板电容器与一电池相连。在电容器的两板间有一带正电的质点处于静止平衡状态。现将电容器两板间的距离增大，则\xzanswer{D}

\fourchoices[answer=D]
{电容变大，质点向上运动}
{电容变大，质点向下运动}
{电容变小，质点保持静止}
{电容变小，质点向下运动}

%% answer: D
%% memo:
\memoanswer{
本题原卷及材料中未提供详解，待补充。
}

\item
%% number: 18
%% paperName: 2005年全国理综Ⅲ第 18 题 6 分
%% typeId: 1
%% type: 单选题
%% chapter: 光学
%% point: 光的电磁说
%% method: 无
%% score: 6
%% degree: 650
%% duplicateId: 0
%% body:
两种单色光由水中射向空气时发生全反射的临界角分别为 $\theta_{1}$、$\theta_{2}$。用 $n_{1}$、$n_{2}$ 分别表示水对两单色光的折射率，$v_{1}$、$v_{2}$ 分别表示两单色光在水中的传播速度\xzanswer{B}

\fourchoices[answer=B]
{$n_{1}<n_{2}$、$v_{1}<v_{2}$}
{$n_{1}<n_{2}$、$v_{1}>v_{2}$}
{$n_{1}>n_{2}$、$v_{1}<v_{2}$}
{$n_{1}>n_{2}$、$v_{1}>v_{2}$}

%% answer: B
%% memo:
\memoanswer{
本题原卷及材料中未提供详解，待补充。
}

\item
%% number: 19
%% paperName: 2005年全国理综Ⅲ第 19 题 6 分
%% typeId: 1
%% type: 单选题
%% chapter: 气体性质
%% point: 热力学第一定律
%% method: 无
%% score: 6
%% degree: 650
%% duplicateId: 0
%% body:
一定质量的气体经历一缓慢的绝热膨胀过程。设气体分子间的势能可忽略，则在此过程中\xzanswer{D}

\fourchoices[answer=D]
{外界对气体做功，气体分子的平均动能增加}
{外界对气体做功，气体分子的平均动能减少}
{气体对外界做功，气体分子的平均动能增加}
{气体对外界做功，气体分子的平均动能减少}

%% answer: D
%% memo:
\memoanswer{
本题原卷及材料中未提供详解，待补充。
}

\item
%% number: 20
%% paperName: 2005年全国理综Ⅲ第 20 题 6 分
%% typeId: 2
%% type: 多选题
%% chapter: 机械波
%% point: 波的图象
%% method: 无
%% score: 6
%% degree: 650
%% duplicateId: 0
%% body:
一列简谐横波在 $x$ 轴上传播，某时刻的波形图如图所示，a、b、c 为三个质元，a 正向上运动。由此可知\xzanswer{AC}

\onepicture[w=6.5cm]{figs/20.png}

\fourchoices[answer=AC]
{该波沿 $x$ 轴正方向传播}
{c 正向上运动}
{该时刻以后，b 比 c 先到达平衡位置}
{该时刻以后，b 比 c 先到达离平衡位置最远处}

%% answer: AC
%% memo:
\memoanswer{
本题原卷及材料中未提供详解，待补充。
}

\item
%% number: 21
%% paperName: 2005年全国理综Ⅲ第 21 题 6 分
%% typeId: 2
%% type: 多选题
%% chapter: 曲线运动
%% point: 人造地球卫星
%% method: 无
%% score: 6
%% degree: 450
%% duplicateId: 0
%% body:
最近，科学家在望远镜中看到太阳系外某一恒星有一行星，并测得它围绕该恒星运行一周所用的时间为 1200 年，它与该恒星的距离为地球到太阳距离的 100 倍。假定该行星绕恒星运行的轨道和地球绕太阳运行的轨道都是圆周，仅利用以上两个数据可以求出的量有\xzanswer{AD}

\fourchoices[answer=AD]
{恒星质量与太阳质量之比}
{恒星密度与太阳密度之比}
{行星质量与地球质量之比}
{行星运行速度与地球公转速度之比}

%% answer: AD
%% memo:
\memoanswer{
本题原卷及材料中未提供详解，待补充。
}

\item
%% number: 22
%% paperName: 2005年全国理综Ⅲ第 22 题 12 分
%% typeId: 4
%% type: 实验题
%% chapter: 电路
%% point: 测电动势与内阻
%% method: 无
%% score: 12
%% degree: 650
%% duplicateId: 0
%% body:
\begin{enumerate}
	\item
	用螺旋测微器测圆柱体的直径时，示数如图所示，此示数为\tkanswer[2.5cm]{$8.116\Umm$}。

	\onepicture[w=5.0cm, align=right]{figs/22a.png}

	\item
	利用图中给定的器材测量电压表 V 的内阻 $R_{V}$。图中 $E$ 为电源（内阻可忽略不计），$R$ 为电阻箱，K 为电键。

	\begin{enumerate}
		\item
		将图中实物连接为测量所用的电路。

		\drawpicanswer{%
		\onepicture[w=5.5cm, align=right]{figs/22b.png}%
		}{%
		\onepicture[w=5.5cm, align=right]{figs/22b_an.png}%
		}

		\item
		写出实验中必须记录的数据（用符号表示），并指出各符号的意义：

		\tkanswer[8cm]{$R_{1}$、$R_{2}$ 是电阻箱的两次读数，$U_{1}$、$U_{2}$ 是相应的电压表的两次读数}。

		\item
		用上述记录的数据表示 $R_{V}$ 的公式为 $R_{V}=$\tkanswer[3cm]{$\frac{U_{2}R_{2}-U_{1}R_{1}}{U_{1}-U_{2}}$}。
	\end{enumerate}
\end{enumerate}

%% answer: 
\jdanswer{
\begin{enumerate}
	\item
	8.116（5 分，在 $8.116\pm0.002$ 范围内都给 5 分）

	\item
	连线如图所示；$R_{V}=\frac{U_{2}R_{2}-U_{1}R_{1}}{U_{1}-U_{2}}$。
\end{enumerate}
}

%% memo:
\memoanswer{
本题原卷及材料中未提供详解，待补充。
}

\item
%% number: 23
%% paperName: 2005年全国理综Ⅲ第 23 题 16 分
%% typeId: 6
%% type: 计算题
%% chapter: 磁场
%% point: 洛伦兹力
%% method: 无
%% score: 16
%% degree: 650
%% duplicateId: 0
%% body:
图中 MN 表示真空室中垂直于纸面的平板，它的一侧有匀强磁场，磁场方向垂直纸面向里，磁感应强度大小为 $B$。一带电粒子从平板上的狭缝 O 处以垂直于平板的初速 $v$ 射入磁场区域，最后到达平板上的 P 点。已知 $B$、$v$ 以及 P 到 O 的距离 $l$。不计重力，求此粒子的电荷 $q$ 与质量 $m$ 之比。

\onepicture[h=4.0cm, align=right]{figs/23.png}

%% answer: 
\jdanswer{$\frac{2v}{Bl}$}

%% memo:
\memoanswer{
粒子初速 $v$ 垂直于磁场，粒子在磁场中受洛伦兹力而做匀速圆周运动，设其半径为 $R$，由洛伦兹力公式和牛顿第二定律，有

$qvB=\frac{mv^{2}}{R}$

因粒子经 O 点时的速度垂直于 OP，故 OP 是直径，$l=2R$，由此得

$\frac{q}{m}=\frac{2v}{Bl}$
}

\item
%% number: 24
%% paperName: 2005年全国理综Ⅲ第 24 题 19 分
%% typeId: 6
%% type: 计算题
%% chapter: 运动和力
%% point: 牛顿第二定律
%% method: 无
%% score: 19
%% degree: 450
%% duplicateId: 0
%% body:
如图所示，在倾角为 $\theta$ 的光滑斜面上有两个用轻质弹簧相连接的物块 A、B。它们的质量分别为 $m_{A}$、$m_{B}$，弹簧的劲度系数为 $k$，C 为一固定挡板。系统处于静止状态。现开始用一恒力 $F$ 沿斜面方向拉物块 A 使之向上运动，求物块 B 刚要离开 C 时物块 A 的加速度 $a$ 和从开始到此时物块 A 的位移 $d$，重力加速度为 $g$。

\onepicture[w=6.0cm, align=right]{figs/24.png}

%% answer: 
\jdanswer{$a=\frac{F-(m_{\mathrm{A}}+m_{\mathrm{B}})g\sin\theta}{m_{\mathrm{A}}}$，$d=\frac{(m_{\mathrm{A}}+m_{\mathrm{B}})g\sin\theta}{k}$}

%% memo:
\memoanswer{
令 $x_{1}$ 表示未加 $F$ 时弹簧的压缩量，由胡克定律和牛顿定律可知

$m_{\mathrm{A}}g\sin\theta=kx_{1}$ ①

令 $x_{2}$ 表示 B 刚要离开 C 时弹簧的伸长量，$a$ 表示此时 A 的加速度，由胡克定律和牛顿定律可知：

$kx_{2}=m_{\mathrm{B}}g\sin\theta$ ②

$F-m_{\mathrm{A}}g\sin\theta-kx_{2}=m_{\mathrm{A}}a$ ③

由②③式可得 $a=\frac{F-(m_{\mathrm{A}}+m_{\mathrm{B}})g\sin\theta}{m_{\mathrm{A}}}$ ④

由题意：$d=x_{1}+x_{2}$ ⑤

由④⑤式可得 $d=\frac{(m_{\mathrm{A}}+m_{\mathrm{B}})g\sin\theta}{k}$
}

\item
%% number: 25
%% paperName: 2005年全国理综Ⅲ第 25 题 20 分
%% typeId: 6
%% type: 计算题
%% chapter: 运动和力
%% point: 动量守恒定律
%% method: 无
%% score: 20
%% degree: 450
%% duplicateId: 0
%% body:
如图所示，一对杂技演员（都视为质点）乘秋千（秋千绳处于水平位置）从 A 点由静止出发绕 O 点下摆，当摆到最低点 B 时，女演员在极短时间内将男演员沿水平方向推出，然后自己刚好能回到高处 A。求男演员落地点 C 与 O 点的水平距离 $s$。已知男演员质量 $m_{1}$ 和女演员质量 $m_{2}$ 之比 $\frac{m_{1}}{m_{2}}=2$，秋千的质量不计，秋千的摆长为 $R$，C 点比 O 点低 $5R$。

\onepicture[h=4.5cm, align=right]{figs/25.png}

%% answer: 
\jdanswer{$s=8R$}

%% memo:
\memoanswer{
设分离前男女演员在秋千最低点 B 的速度为 $v_{0}$，由机械能守恒定律

$(m_{1}+m_{2})gR=\frac{1}{2}(m_{1}+m_{2})v_{0}^{2}$

设刚分离时男演员速度的大小为 $v_{1}$，方向与 $v_{0}$ 相同；女演员速度的大小为 $v_{2}$，方向与 $v_{0}$ 相反，由动量守恒：

$(m_{1}+m_{2})v_{0}=m_{1}v_{1}-m_{2}v_{2}$

分离后，男演员做平抛运动，设男演员从被推出到落在 C 点所需的时间为 $t$，根据题给条件，由运动学规律：$4R=\frac{1}{2}gt^{2}$，$s=v_{1}t$

根据题给条件，女演员刚好回到 A 点，由机械能守恒定律，$m_{2}gR=\frac{1}{2}m_{2}v_{2}^{2}$

已知 $\frac{m_{1}}{m_{2}}=2$，由以上各式可得 $s=8R$
}

\end{enumerate}

\end{document}
